\section{Vận hành, artifact và FID}

\subsection{Sampling Stage 2 trên đường JAX}

\subsubsection{Sampling nhanh}

\begin{lstlisting}[style=rae]
python3 src_jax/sample.py \
  --config <sample_config> \
  --seed 42 \
  --class-labels 207,360 \
  --output sample_jax.png
\end{lstlisting}

\subsubsection{Sampling phân tán}

\begin{lstlisting}[style=rae]
python3 src_jax/sample_ddp.py \
  --config <sample_config> \
  --sample-dir samples_jax \
  --num-samples 50000 \
  --per-proc-batch-size 4 \
  --label-sampling equal \
  --fid-ref /path/to/reference_stats.npz
\end{lstlisting}

Artifact đầu ra:

\begin{itemize}[leftmargin=1.5em]
  \item các file PNG theo index;
  \item các shard \texttt{.npz} theo rank nếu bật \texttt{--save-npz};
  \item giá trị FID in ra terminal nếu cung cấp \texttt{--fid-ref}.
\end{itemize}

\subsection{Build reference statistics cho FID}

\begin{lstlisting}[style=rae]
python src/build_fid_stats.py \
  --input /path/to/reference_images \
  --output /path/to/reference_stats.npz \
  --device cpu \
  --batch-size 64 \
  --num-workers 32 \
  --num-threads 96
\end{lstlisting}

Đây là đường \texttt{torch-fidelity} chung cho nhánh XLA/PyTorch. Nếu cần
\texttt{fid\_ref} cho online FID ở JAX, nên build bằng detector backend-native:

\begin{lstlisting}[style=rae]
python3 src_jax/build_fid_stats.py \
  --input /path/to/reference_images \
  --output /path/to/reference_stats.pkl \
  --batch-size 64 \
  --num-workers 8
\end{lstlisting}

\subsection{Artifact của Stage 2 training}

\subsubsection{Đường XLA}

Một experiment trong \texttt{results/} thường có:

\begin{itemize}[leftmargin=1.5em]
  \item \texttt{log.txt},
  \item \texttt{checkpoints/<step>.pt},
  \item \texttt{fid\_eval/step\_<step>/...} nếu bật online FID.
\end{itemize}

\subsubsection{Đường JAX}

Một run trong \texttt{results\_jax/} thường có:

\begin{itemize}[leftmargin=1.5em]
  \item \texttt{jax\_adapter\_config.json},
  \item thư mục checkpoint Orbax do backend quản lý,
  \item media và scalar trên wandb nếu bật \texttt{--wandb},
  \item toàn bộ workdir có thể được đẩy lên Hugging Face nếu truyền
  \texttt{--hf-repo-id}.
\end{itemize}

\subsection{Đẩy model lên Hugging Face}

\begin{lstlisting}[style=rae]
python3 src_jax/push_hf.py \
  --path results_jax/<run_name> \
  --repo-id <hf_user_or_org>/<repo_name>
\end{lstlisting}

Hoặc gắn thẳng vào lệnh train:

\begin{lstlisting}[style=rae]
python3 src_jax/train.py \
  --config <config> \
  --data-path <train_root> \
  --hf-repo-id <hf_user_or_org>/<repo_name>
\end{lstlisting}

\subsection{Các lưu ý vận hành quan trọng}

\begin{itemize}[leftmargin=1.5em]
  \item Đường JAX cố ý giữ nguyên schema YAML của repo để giảm chi phí bảo trì.
  \item \texttt{stage\_2.ckpt} trên đường JAX có thể là PyTorch hoặc Orbax, nhưng
  cách load sẽ khác nhau.
  \item Đường JAX giờ cũng có validation loss định kỳ qua block \texttt{eval},
  ngoài online FID.
  \item FID trên đường JAX nên dùng reference stats build bởi
  \texttt{src\_jax/build\_fid\_stats.py} để cùng detector với backend; nếu cần,
  adapter vẫn chuyển \texttt{.npz} cũ sang pickle nội bộ.
  \item Online FID ở JAX mặc định chấm EMA. Nếu cần đối chiếu với model hiện
  thời, bật \texttt{eval.fid\_eval\_model=true} để log thêm metric
  \texttt{FID-4K/model (cfg=...)} bên cạnh series EMA mặc định.
  \item Nếu chỉ cần inference Stage 1 ở JAX, dùng \texttt{src\_jax/stage1\_sample.py}.
  \item Nếu cần stat Stage 1 hoặc reconstruction cả thư mục ở JAX, dùng
  \texttt{src\_jax/build\_stage1\_stats.py} và \texttt{src\_jax/reconstruct\_folder.py}.
  Các lệnh này có thể dùng YAML chỉ có \texttt{stage\_1}, không cần
  \texttt{stage\_2.target}.
  \item Nếu chạy trên Kaggle với CelebA-HQ, bắt đầu từ notebook
  \texttt{vaes-jax-celebahq-kaggle.ipynb}. Bản này vẫn dùng
  \texttt{src\_jax/export\_celebahq\_hf.py}, nhưng đổi Stage 1 sang
  \texttt{stage1.StabilityVAE}, đổi Stage 2 sang single-tower
  \texttt{SiT-B}, và bỏ hẳn bước bootstrap identity stat cùng pass
  \texttt{build\_stage1\_stats.py} bắt buộc.
  \item Nếu phần cứng là Kaggle \texttt{TPU v5e-8} và cần flow VAE đó trên TPU,
  dùng \texttt{vaes-jax-celebahq-kaggle-tpuv5e8-sitb.ipynb}. Notebook này giữ
  toàn bộ fix Kaggle/JAX hiện có, bật \texttt{training.random\_flip=true},
  viết config Stage 2 thành
  \texttt{CelebAHQ256\_SiT-B\_StabilityVAE\_jax\_tpuv5e8.yaml}, và dùng shape
  single-tower \texttt{hidden\_size=768}, \texttt{depth=12},
  \texttt{num\_heads=12}.
  \item Nếu đã có run Orbax của flow \texttt{StabilityVAE + SiT-B} và cần
  resume, dùng notebook
  \texttt{vaes-jax-celebahq-kaggle-tpuv5e8-sitb-resume.ipynb}. Notebook này
  dùng cùng pattern resume-only của bản DH, nhưng tự tìm thư mục
  \texttt{CelebAHQ256\_SiT-B\_StabilityVAE\_jax\_tpuv5e8-*} mới nhất.
  \item Nếu cần huấn luyện Stage 1 kiểu GAN đầy đủ, hiện vẫn phải xem đây là
  phần chưa được JAX hóa trong branch này.
\end{itemize}

\subsection{Checklist thực tế trước khi chạy dài}

\begin{enumerate}[leftmargin=1.8em]
  \item Xác nhận \texttt{stage\_1} và \texttt{stage\_2} khớp latent shape.
  \item Xác nhận dữ liệu train và val đúng dạng \texttt{ImageFolder}.
  \item Với flow RAE và dataset mới, build lại \texttt{stage1\_stat.pt} thay vì
  tái dùng stat của ImageNet nếu muốn latent normalization khớp dữ liệu hơn.
  \item Build \texttt{fid\_ref} trước nếu muốn FID ổn định.
  \item Export đủ biến wandb trước khi bật \texttt{--wandb}.
  \item Nếu muốn upload lên Hugging Face, xác nhận token nằm trong biến môi
  trường được chỉ định bởi \texttt{--hf-token-env}.
  \item Với JAX sampling số lượng lớn, chọn \texttt{--per-proc-batch-size} đủ an
  toàn theo bộ nhớ thiết bị.
\end{enumerate}
